\subsection{Primordial Saturation and Direct Collapse}
  \label{subsec:primordial_saturation}

  The redshift dependence discussed above acquires particular significance
  in the primordial regime.
  As introduced in Section~\ref{subsec:operational_saturation_scale},
  the operational saturation scale
  \begin{equation}
    a_\star(z) \sim \kappa \frac{c_{\chi}^2}{\ell(z)}
  \end{equation}
  depends inversely on the descriptive resolution~$\ell(z)$.
  At very high redshift, the projected state remains spectrally constrained,
  so that $\ell(z)$ cannot decrease arbitrarily while preserving smooth
  quasi-injective updates.

  Consider a collapsing primordial gas cloud of mass~$M$ and radius~$R$.
  The characteristic update demand is set by the free-fall timescale,
  \begin{equation}
    D(R) \sim \frac{1}{t_{\mathrm ff}} \sim \sqrt{\frac{GM}{R^3}}.
  \end{equation}
  Fragmentation into stellar units requires that the characteristic
  fragmentation scale $\lambda_{\mathrm frag}$ remain larger than the minimal
  projectable scale~$\ell(z)$.
  When $\lambda_{\mathrm frag} < \ell(z)$, the projection cannot resolve a
  multi-body stellar configuration.

  In this regime, admissibility can only be restored through a saturated
  reconfiguration channel.
  Direct collapse into a compact seed then becomes the structurally
  available outcome.
  This yields an implicit redshift-dependent critical scale
  $M_{\mathrm crit}(z)$ defined by the condition that the required
  fragmentation scale equals~$\ell(z)$.
  Massive halos above this threshold naturally enter a direct-collapse
  regime without requiring additional exotic physics.
